\documentclass[10pt,a4paper]{amsart}
\usepackage[margin=19mm]{geometry}
\usepackage{amsmath,amssymb,mathtools}
\usepackage[scr=rsfs,cal=euler]{mathalfa}
\usepackage{xfrac}
\usepackage[unicode,hidelinks,bookmarks=false]{hyperref}
\newcommand{\ie}{i.e.\ }
\newcommand{\cf}{cf.\ }
\newcommand{\resp}{respectively\ }
\newcommand{\Subsection}{\subsection}
\providecommand{\nobreakdash}{\nobreak\hbox{-}}
\setlength{\emergencystretch}{2em}
\multlinegap0pt
\usepackage{bm}
\usepackage{bbm}
\usepackage{dsfont}
\usepackage{xspace}
\usepackage{cancel}
\usepackage{mathtools}
\usepackage{amsthm}
\makeatletter
\@ifundefined{theorem}{\newtheorem{theorem}{Theorem}}{}
\makeatother
\newcommand{\A}{\mathcal{A}} 
\newcommand{\B}{\mathcal{B}} 
\title[High-Order Asymptotic-Preserving Schemes for Kinetic Equatio]{High-Order Asymptotic-Preserving Schemes for Kinetic Equations from Rarefied to Incompressible Regimes}
\author{Giacomo Dimarco, Axel Klar, Theresa K\"ofler, Lorenzo Pareschi, Tiwari Sudarshan}
\date{Source: arXiv:2606.28040v1 (CC BY 4.0)}
\begin{document}
\maketitle

\noindent\emph{Source and licence.} High-Order Asymptotic-Preserving Schemes for Kinetic Equations from Rarefied to Incompressible Regimes. Version arXiv:2606.28040v1 (CC BY 4.0), distributed under the Creative Commons Attribution 4.0 International licence, as recorded in the arXiv licence metadata for this exact version and shown on the abstract page https://arxiv.org/abs/2606.28040v1. Full rights notice: licenses/arxiv-2606.28040-CC-BY-4.0.txt.\par\smallskip

\noindent\textit{Editorial scope.} Counted unit: the Theorem whose source label is \texttt{ThmFullBGKTypeA} in the article (source0.tex lines 999--1001, source label \texttt{ThmFullBGKTypeA}), delivered together with the article's own complete original proof of it (source0.tex lines 1002--1192, one \texttt{proof} environment of 191 lines). No mathematics is written, repaired or re-derived: statement and proof are the article's own text line for line. Required context retained uncounted: InternalStages (lines 989--997). Every definition, display and label of those lines is retained unchanged. Omitted from this package and not redistributed: 1--988, 998--998, 1193--2512. Nothing inside a retained excerpt is omitted or altered: every retained body is delivered exactly as the article's lines are written, so no omission receipt of any kind accompanies this package. Citations: the retained excerpts carry no citation command at all, so no bibliography entry of the article is transcribed and no reference is redistributed. Reference adaptation: none was needed and none was made; every reference inside the delivered unit resolves inside it, so no unresolved reference or citation is produced. Printed slips kept literal: no printed word, symbol, hypothesis or number of the counted statement or of its proof was altered. Layout and class: the article's own class is \texttt{siamart251216} with packages including babel, framed, multirow, hyperref, algpseudocode, amssymb, amsmath, bm, amsfonts, algorithm, algcompatible, dsfont, latexsym, subcaption; the wrapper is the lane's pinned \texttt{amsart} editorial wrapper at 10pt on A4, which restates the theorem-like environments the retained text needs and adds the article's own macro definitions that the retained text needs (\texttt{\textbackslash A}, \texttt{\textbackslash B}), with extra wrapper packages (bm, bbm, dsfont, xspace, cancel, mathtools). The delivered numbering is the unit's own. Assets: no retained span contains a figure environment, a graphics-inclusion command, a PDF-page inclusion, a table or a TikZ picture; no image file is copied into this package, the package declares no asset dependency and no image file is redistributed with it. Licence: version arXiv:2606.28040v1 of this article is distributed under the Creative Commons Attribution 4.0 International licence, as recorded in the arXiv licence metadata for this exact version and shown on the abstract page https://arxiv.org/abs/2606.28040v1. A full rights notice is delivered at licenses/arxiv-2606.28040-CC-BY-4.0.txt. Characters: every retained excerpt is pure ASCII, so the delivered engine, XeLaTeX with its default Latin Modern text font, reports no missing character.
\begin{align}
\label{InternalStages}
\rho^{(i)} &= \rho^n - \frac{\Delta t}{\varepsilon} \sum_{j = 1}^i a_{ij} \nabla_x \cdot u^{(j)}, \nonumber \\
u^{(i)} &= u^n - \Delta t  \sum_{j = 1}^i \left( a_{ij} \nabla_x \cdot \langle \mathcal{A}(v) g^{(j)} \rangle  + \frac{a_{ij}}{\varepsilon} \nabla_x (T^{(j)} + \rho^{(j)})\right), \nonumber \\
T^{(i)} &= T^n - \Delta t  \sum_{j = 1}^i \left( \frac{a_{ij}}{d} \nabla_x \cdot \langle \mathcal{B}(v) g^{(j)} \rangle  + \frac{2 a_{ij}}{d \varepsilon} \nabla_x \cdot u^{(j)}\right), \nonumber \\
p^{(i)} &= p^n - \Delta t  \sum_{j = 1}^i \left( \frac{a_{ij}}{d} \nabla_x \cdot \langle \mathcal{B}(v) g^{(j)} \rangle  + \frac{d+2}{d \varepsilon} a_{ij} \nabla_x \cdot u^{(j)}\right), \\
g^{(i)} &= g^n - \frac{\Delta t}{\varepsilon^2} \sum_{j=1}^{i-1} \tilde{a}_{ij} \Bigl( \nabla_x \cdot \left[ \mathcal{A}(v)u^{(j)} + \tfrac{\mathcal{B}(v)}{2}T^{(j)} + \varepsilon (I-P) (vg^{(j)})\right] \nonumber \\
&\quad - Q(u^{(j)},\rho^{(j)},T^{(j)}) \Bigr) - \frac{\Delta t}{\varepsilon^2 \tau} \sum_{j=1}^i a_{ij} g^{(j)}. \nonumber
\end{align}

\begin{theorem}\label{ThmFullBGKTypeA}
	If the IMEX method is of type A and satisfies the GSA property, then in the fluid dynamic limit $\varepsilon \rightarrow 0$ the IMEX RK scheme \eqref{InternalStages} becomes a consistent discretization for the thermal  incompressible Navier Stokes equations.
\end{theorem}

\begin{proof}
Rewriting the last equation of \eqref{InternalStages} gives
%\begin{equation}
%\begin{split} \label{InternalStagesForg}
%g^{(i)} &= \frac{\varepsilon^2 \tau}{\varepsilon^2 \tau + a_{ii} \Delta t}g^n - \frac{\Delta t \tau}{\varepsilon^2 \tau + a_{ii} \Delta t} \sum_{j=1}^{i-1} \tilde{a}_{ij}\left( \nabla_x \cdot \left[ \mathcal{A}(v)u^{(j)} + \frac{\mathcal{B}(v)}{2}T^{(j)} + \varepsilon (I-P) (vg^{(j)})\right] - Q^{(j)} \right) \\&- \frac{\Delta t}{\varepsilon^2 \tau + a_{ii} \Delta t} \sum_{j=1}^{i-1} a_{ij} g^{(j)},
%\end{split}
%\end{equation}
\begin{align}
	\label{InternalStagesForg}
&g^{(i)} = \, \frac{\varepsilon^2 \tau}{\varepsilon^2 \tau + a_{ii} \Delta t} g^n \nonumber \\
&- \frac{\Delta t \tau}{\varepsilon^2 \tau + a_{ii} \Delta t} \sum_{j=1}^{i-1} \tilde{a}_{ij} \Biggl( \nabla_x \cdot \left[ \mathcal{A}(v) u^{(j)} + \frac{\mathcal{B}(v)}{2} T^{(j)} + \varepsilon (I-P)(vg^{(j)}) \right] - Q^{(j)} \Biggr) \\
&- \frac{\Delta t}{\varepsilon^2 \tau + a_{ii} \Delta t} \sum_{j=1}^{i-1} a_{ij} g^{(j)},\nonumber
\end{align}
where $Q^{(j)} = Q(u^{(j)},\rho^{(j)},T^{(j)})$ for readability. Then, by inserting  this %\eqref{InternalStagesForg} 
into the second equation of \eqref{InternalStages}, using orthogonality of $\mathcal{A}(v)$ and $\mathcal{B}(v)$  and taking the limit $\varepsilon \to 0$ we obtain
%\begin{equation}
%\begin{split}
%u^{(i)} &= u^n - \Delta t  \sum_{j = 1}^{i-1} \left( a_{ij} \nabla_x \cdot \langle \mathcal{A}(v) g^{(j)} \rangle \right)
%+ \Delta t \tau  \sum_{j=1}^{i-1} \tilde{a}_{ij}  \nabla_x \cdot \left( \langle \mathcal{A}(v) \otimes  \mathcal{A}(v) \rangle : \nabla_x u^{(j)} - \langle \mathcal{A}(v) Q^{(j)} \rangle\right) \\& + \Delta t  \sum_{j = 1}^{i-1} \left( a_{ij} \nabla_x \cdot \langle \mathcal{A}(v) g^{(j)} \rangle \right)
%- \frac{\Delta t}{\varepsilon}  \sum_{j = 1}^i  a_{ij} \nabla_x (T^{(j)} + \rho^{(j)}).
%\end{split}
%\end{equation}
\begin{equation}
\begin{split}
u^{(i)} = \, & u^n - \Delta t  \sum_{j = 1}^{i-1} a_{ij} \nabla_x \cdot \langle \mathcal{A}(v) g^{(j)} \rangle \\
& + \Delta t \tau  \sum_{j=1}^{i-1} \tilde{a}_{ij} \nabla_x \cdot \left( \langle \mathcal{A}(v) \otimes \mathcal{A}(v) \rangle : \nabla_x u^{(j)} - \langle \mathcal{A}(v) Q^{(j)} \rangle \right) \\
& + \Delta t  \sum_{j = 1}^{i-1} a_{ij} \nabla_x \cdot \langle \mathcal{A}(v) g^{(j)} \rangle - \frac{\Delta t}{\varepsilon} \sum_{j = 1}^i a_{ij} \nabla_x (T^{(j)} + \rho^{(j)}).
\end{split}
\end{equation}
Rescaling $p$ as  $\varepsilon \tilde p $ and, for simplicity, dropping  $\widetilde{~}$ ~in what follows, we obtain
\begin{equation}
\begin{split} \label{InternalStagesforu}
u^{(i)} &= u^n 
+ \Delta t \tau  \sum_{j=1}^{i-1} \tilde{a}_{ij}  \nabla_x \cdot \left( \langle \mathcal{A}(v) \otimes  \mathcal{A}(v) \rangle : \nabla_x u^{(j)} - \langle \mathcal{A}(v) Q^{(j)} \rangle\right) 
- \Delta t  \sum_{j = 1}^i  a_{ij} \nabla_x p^{(j)}.
\end{split}
\end{equation}
We now consider the stage equations for $p$, that can be directly obtained from the equations for $\rho$ and $T$.
\begin{equation}
\varepsilon p^{(i)} = \varepsilon p^n  - \Delta t  \sum_{j = 1}^i \left( \frac{1}{d} a_{ij} \nabla_x \cdot \langle \mathcal{B}(v) g^{(j)} \rangle  + \frac{2+d}{d \varepsilon} a_{ij} \nabla_x \cdot u^{(j)}\right) 
\end{equation}
Substituting $u^{(i)}$ gives
%\begin{equation}\begin{split}
%p^{(i)} &= p^n  - \Delta t  \sum_{j = 1}^i \left( \frac{1}{d \varepsilon} a_{ij} \nabla_x \cdot \langle \mathcal{B}(v) g^{(j)} \rangle \right) 
%- \Delta t  \sum_{j = 1}^{i-1} \left( \frac{2+d}{d \varepsilon^2} a_{ij} \nabla_x \cdot u^{(j)}\right)  -  \frac{2+d}{d}  \frac{\Delta t}{ \varepsilon^2} a_ii \nabla_x \cdot u^{(i)}_* \\&+ \frac{2+d}{d}  \frac{\Delta t^2}{ \varepsilon^2} a_{ii}^2 \Delta_x p^{(i)},
%\end{split}
%\end{equation}
\begin{equation}
\label{eq:pressure_stage}
\begin{split}
p^{(i)} = \, & p^n - \frac{\Delta t}{d \varepsilon} \sum_{j = 1}^i a_{ij} \nabla_x \cdot \langle \mathcal{B}(v) g^{(j)} \rangle 
- \frac{(d+2)\Delta t}{d \varepsilon^2} \sum_{j = 1}^{i-1} a_{ij} \nabla_x \cdot u^{(j)} \\
& - \frac{d+2}{d} \frac{\Delta t}{\varepsilon^2} a_{ii} \nabla_x \cdot u^{(i)}_* 
+ \frac{d+2}{d} \frac{\Delta t^2}{\varepsilon^2} a_{ii}^2 \Delta_x p^{(i)}.
\end{split}
\end{equation}
where
\begin{equation}
u^{(i)}_* = u^n - \Delta t  \sum_{j = 1}^i \left( a_{ij} \nabla_x \cdot \langle \mathcal{A}(v) g^{(j)} \rangle \right) - \Delta t  \sum_{j = 1}^{i-1} \left( a_{ij} \nabla_x p^{(j)} \right).
\end{equation}
Rewriting gives
%\begin{equation}\begin{split}\label{InternalStagesforp}
%\frac{2+d}{d}   \Delta_x p^{(i)} - \frac{ \varepsilon^2}{\Delta t^2 a_{ii}^2} p^{(i)} &= -\frac{ \varepsilon^2}{\Delta t^2 a_{ii}^2} p^n  + \frac{ \varepsilon^2}{\Delta t a_{ii}^2}  \sum_{j = 1}^i \left( \frac{1}{d \varepsilon} a_{ij} \nabla_x \cdot \langle \mathcal{B}(v) g^{(j)} \rangle \right) 
%+ \frac{ 1}{\Delta t a_{ii}^2}  \sum_{j = 1}^{i-1} \left( \frac{2+d}{d } a_{ij} \nabla_x \cdot u^{(j)}\right)  \\&+  \frac{ 1}{\Delta t a_{ii}} \frac{2+d}{d} \nabla_x \cdot u^{(i)}_*.
%\end{split}
%\end{equation}
\begin{equation}
\label{InternalStagesforp}
\begin{split}
\frac{d+2}{d} a_{ii}^2 \Delta_x p^{(i)} - \frac{\varepsilon^2}{\Delta t^2 } p^{(i)} = \, & -\frac{\varepsilon^2}{\Delta t^2 } p^n + \frac{\varepsilon}{d \Delta t } \sum_{j = 1}^i a_{ij} \nabla_x \cdot \langle \mathcal{B}(v) g^{(j)} \rangle \\
& + \frac{d+2}{d \Delta t } \sum_{j = 1}^{i-1} a_{ij} \nabla_x \cdot u^{(j)} + \frac{(d+2)a_{ii}}{d \Delta t } \nabla_x \cdot u^{(i)}_*.
\end{split}
\end{equation}
Taking the divergence of the second equation of \eqref{InternalStages} gives for the first stage
\begin{equation} \label{SecondProofReference}
\nabla_x \cdot u^{(1)} = \nabla_x \cdot u^n - \Delta t  a_{11} \nabla_x \cdot \nabla_x \cdot \langle \mathcal{A}(v) g^{(1)} \rangle
- \Delta t a_{11} \Delta_x p^{(1)}.
\end{equation}
In the limit $\varepsilon \to 0$, we get for the first stage of \eqref{InternalStagesforp}
\begin{equation}
\Delta t a_{11} \Delta_x p^{(1)} =  \nabla_x \cdot u^n - \Delta t  a_{11} \nabla_x \cdot \nabla_x \cdot \langle \mathcal{A}(v) g^{(1)} \rangle
\end{equation}
and therefore $\nabla_x \cdot u^{(1)} = 0$.
For the inductive step $(i-1) \mapsto i$, the second equation of \eqref{InternalStages} gives
\begin{equation}
\begin{split}
\nabla_x \cdot u^{(i)} &= \nabla_x \cdot u^n 
- \Delta t  \sum_{j = 1}^i \left( a_{ij} \nabla_x \cdot \nabla_x \cdot \langle \mathcal{A}(v) g^{(j)} \rangle \right) - \Delta t  \sum_{j = 1}^{i-1} \left( a_{ij} \Delta_x p^{(j)} \right) 
- \Delta t  a_{ii} \Delta_x p^{(i)}.
\end{split}
\end{equation}
Using \eqref{InternalStagesforp} and the induction hypothesis $\nabla_x \cdot u^{(i-1)} = 0$ gives in the limit 
\begin{equation}
\begin{split}
\Delta t a_{ii} \Delta_x p^{(i)}  &= \nabla_x \cdot  u^n - \Delta t  \sum_{j = 1}^i \left( a_{ij} \nabla_x \cdot \nabla_x \cdot \langle \mathcal{A}(v) g^{(j)} \rangle \right) - \Delta t  \sum_{j = 1}^{i-1} \left( a_{ij} \Delta_x p^{(j)} \right).
\end{split}
\end{equation}
and therefore we obtain $\nabla_x \cdot u^{(i)} = 0$ in the incompressible limit.
By using the following properties
\begin{equation*}\begin{split}
\nabla_x \cdot \langle \mathcal{A}(v) Q \rangle &= \frac{1}{\tau}\left( \nabla_x \cdot \mathcal{A}(u)\right) = \frac{1}{\tau} \left((\nabla_x \cdot u) u + (u \cdot \nabla_x) u - \nabla_x \frac{|u|^2}{d} \right)\\
\nabla_x \cdot \langle \mathcal{A}(v) \otimes \mathcal{A}(v) \rangle : \nabla_x u &= \Delta_x u + \nabla_x(\nabla_x \cdot u) - \frac{2}{d} \nabla_x (\nabla_x \cdot u)
\end{split}
\end{equation*}
as well as the fact that all internal stages of $u$ are divergence-free in the incompressible limit, we obtain
\begin{equation}
\begin{split} \label{FinalForu_}
u^{(i)} &= u^n 
+ \Delta t  \sum_{j=1}^{i-1} \tilde{a}_{ij}  \left( \tau \Delta_x u^{(j)} -  (u^{(j)} \cdot \nabla_x ) u^{(j)}\right) 
+ \Delta t  \sum_{j=1}^{i-1} \tilde{a}_{ij} \nabla_x \frac{|u^{(j)}|^2}{d} 
- \Delta t  \sum_{j = 1}^i  a_{ij} \nabla_x p^{(j)}.
\end{split}
\end{equation}
Using Bernoulli's relation finally gives 
\begin{equation}
\begin{split} \label{FinalForu}
u^{(i)} &= u^n 
+ \Delta t  \sum_{j=1}^{i-1} \tilde{a}_{ij}  \left( \tau \Delta_x u^{(j)} -  (u^{(j)} \cdot \nabla_x ) u^{(j)}\right) 
- \Delta t  \sum_{j = 1}^i  a_{ij} \nabla_x p^{(j)},
\end{split}
\end{equation}
where $p$ now denotes the static pressure.
We now consider the stage equations for $p$. Taking the limit $\varepsilon \to 0$ and using the divergence-free condition for all internal stages of $u$ as well as Bernoulli's relation, we get
%\begin{equation}\begin{split} \label{FinalForp_}
%  \sum_{j = 1}^{i} \left( a_{ij} \Delta_x p^{(j)} \right)  &=  \frac{1}{\Delta t}\nabla_x \cdot \left( u^n + \Delta t  \sum_{j=1}^{i-1} \tilde{a}_{ij}  \nabla_x \cdot \left( \tau \Delta_x u^{(j)} -  (\nabla_x \cdot u^{(j)}) u^{(j)}\right)  \right) +  \sum_{j=1}^{i-1} \tilde{a}_{ij} \Delta_x \frac{|u^{(j)}|^2}{d}.
%\end{split}
%\end{equation}
%Using the incompressible and non-dimensional form of Bernoulli's relation, we finally obtain
\begin{equation}\begin{split} \label{FinalForp}
  \sum_{j = 1}^{i} \left( a_{ij} \Delta_x p^{(j)} \right)  &=  \frac{1}{\Delta t}\nabla_x \cdot \left( u^n + \Delta t  \sum_{j=1}^{i-1} \tilde{a}_{ij}  
 \left( \tau \Delta_x u^{(j)} -  (u^{(j)} \cdot  \nabla_x ) u^{(j)}\right)  \right).
\end{split}
\end{equation} 
Finally, we consider the stage equations for $T$.
\begin{equation}
		T^{(i)} = T_*^{(i)}- \Delta t  \sum_{j = 1}^i \left(  \frac{2}{d \varepsilon} a_{ij} \nabla_x \cdot u^{(j)}\right) 
\end{equation}
with
\begin{equation}
	\label{Tstar}
	T_*^{(i)} = T^n - \Delta t  \sum_{j = 1}^i \left( \frac{1}{d} a_{ij} \nabla_x \cdot \langle \mathcal{B}(v) g^{(j)} \rangle  \right) .
\end{equation}
 In the limit $\varepsilon \to 0$ by inserting \eqref{InternalStagesForg} into \eqref{Tstar} and using orthogonality of $\A(v)$ and $\B(v)$, we obtain
%\begin{equation}\begin{split}
%T_*^{(i)} &= T^n - \Delta t  \sum_{j = 1}^{i-1} \left( \frac{1}{d} a_{ij} \nabla_x \cdot \langle \mathcal{B}(v) g^{(j)} \rangle  \right) + \frac{\Delta t}{d}  \sum_{j=1}^{i-1} \tilde{a}_{ij} \nabla_x \cdot \left( \frac{\tau}{2} \langle \mathcal{B}(v) \otimes \mathcal{B}(v) \rangle \nabla_x T^{(j)} -  (d+2)(u^{(j)}T^{(j)}) \right) \\&+ \Delta t  \sum_{j = 1}^{i-1} \left( \frac{1}{d} a_{ij} \nabla_x \cdot \langle \mathcal{B}(v) g^{(j)} \rangle  \right).
%\end{split}
%\end{equation}
\begin{equation}
\label{eq:T_star_stage}
\begin{split}
T_*^{(i)} = \, & T^n - \frac{\Delta t}{d} \sum_{j = 1}^{i-1} a_{ij} \nabla_x \cdot \langle \mathcal{B}(v) g^{(j)} \rangle \\
& + \frac{\Delta t}{d} \sum_{j=1}^{i-1} \tilde{a}_{ij} \nabla_x \cdot \left[ \frac{\tau}{2} \langle \mathcal{B}(v) \otimes \mathcal{B}(v) \rangle \nabla_x T^{(j)} - (d+2)(u^{(j)}T^{(j)}) \right] \\
& + \frac{\Delta t}{d} \sum_{j = 1}^{i-1} a_{ij} \nabla_x \cdot \langle \mathcal{B}(v) g^{(j)} \rangle.
\end{split}
\end{equation}
By using the following properties 
\begin{equation}
\begin{split}
\langle \mathcal{B}(v) \otimes \mathcal{B}(v) \rangle = 2(d+2)I , \;\; 
\langle \mathcal{B}(v) Q \rangle = \frac{d+2}{\tau} (uT)
\end{split}
\end{equation}
we get
\begin{equation}\begin{split} 
T_*^{(i)} &= T^n + \Delta t \frac{d+2}{d} \sum_{j=1}^{i-1} \tilde{a}_{ij} \left( \tau \Delta_x T^{(j)} - \nabla_x \cdot (u^{(j)}T^{(j)}) \right).
\end{split}
\end{equation}
Moreover, 
for $\epsilon \rightarrow 0$ we have 
%\begin{eqnarray*}
%	T^{(i)}-T^n =	T^{(i)}- T_*^{(i)} +T_*^{(i)}- T^n = (-\frac{2}{d+2}+1)(T_*^{(i)} -T^n )
%	=  \Delta t\sum_{j=1}^{i-1} \tilde{a}_{ij} [  \tau \Delta_x T^{(j)} - \nabla_x \cdot (u^{(j)}T^{(j)})]
%\end{eqnarray*}
\begin{align*}
T^{(i)} - T^n &= T^{(i)} - T_*^{(i)} + T_*^{(i)} - T^n \\
&= \left( -\frac{2}{d+2} + 1 \right) (T_*^{(i)} - T^n) \\
&= \Delta t \sum_{j=1}^{i-1} \tilde{a}_{ij} \left[ \tau \Delta_x T^{(j)} - \nabla_x \cdot (u^{(j)} T^{(j)}) \right].
\end{align*}

due to
\begin{eqnarray*}
	T^{(i)}-T_*^{(i)} 	
	= -\frac{2}{d+2}(T_*^{(i)}  -T^n )
\end{eqnarray*}
as previously and, altogether
\begin{equation}\begin{split} \label{FinalForT}
		T^{(i)} &= T^n + \Delta t\sum_{j=1}^{i-1} \tilde{a}_{ij} \left( \tau \Delta_x T^{(j)} - \nabla_x \cdot (u^{(j)}T^{(j)}) \right).
	\end{split}
\end{equation}
Since we are considering a GSA IMEX Runge Kutta scheme, the last internal stages of \eqref{FinalForu}, \eqref{FinalForp} and \eqref{FinalForT} form a consistent high order Implicit-Explicit time discretization of the limit equations.
\end{proof}

\end{document}
